\begin{afterword}
\summaryhead

The post retells how the author, at eighteen, held a theory that had experimental consequences and was professed testable. By the traditional ideal of Science, that was all the virtue required. The ideal, the author says, gives out gold stars too readily. It leaves the choice of which hypothesis to test to the individual, and asks only that everyone accept clear experimental results. That is enough for science to advance, if slowly.

Probability theory is stricter. Given a base rate, a hit rate and a false-positive rate, the chance that a woman with a positive mammogram has breast cancer is exactly 7.5\%, with no room for whim. A standard like that could have told the young author in advance that the hypothesis was a poor choice. It is much harder to use than Science, since one error can lead anywhere, and Science has good reason not to trust individuals with it. But someone who does not want to waste ten years must attempt it, and not knowing the priors does not make them a free choice. The post ends by quoting Nick Tarleton: the problem is a private standard as lax as the social one.

\responsehead

The post's closing distinction is clear and useful. A lax social standard, which lets anyone test an odd hypothesis without being called a bad person, is a poor standard for what one privately believes. The post credits Nick Tarleton with that wording, and its mammography arithmetic is correct. Its two main contrasts, between Science and Bayes, are drawn less carefully.

The first contrast is that Science gives no verdict, before the test, on which hypothesis is best to test. The post describes this ideal in its own words and quotes no one who preaches it. The tradition's best-known modern philosopher, Karl Popper, did rank hypotheses before testing. Popper preferred theories with more content, and so, in the Stanford Encyclopedia's summary, held that ``the more improbable a theory is the better it is scientifically.'' So both standards give a criterion before the test. They rank by different things: Popper by boldness, the Bayesian by probability. The post does not mention Popper's criterion.

The second contrast is Bayes's ``exactly right probability estimate'' with ``no flex or room for whims.'' In the mammography example that holds, because the prior is given. In the case the post cares about, where no one gives the prior, it holds only on one side of a dispute among Bayesians. Objective Bayesians such as Jaynes, whom the post names as a source, add norms that constrain priors. Subjective Bayesians hold that every coherent prior is permitted. The post does not mention the dispute. Its argument against free priors, that whim cannot make an estimate accurate, answers a claim about accuracy. The subjective school's claim is about what reason permits.

The claim that Bayes is ``the law'' behind every statistical tool is the thesis of ``Beautiful Probability,'' whose notes assess it.

In short: a clear statement that a lax public standard makes a poor private one, set against an account of Science that leaves out Popper's criterion for choosing hypotheses, and an account of Bayes that belongs to one of its schools.
\end{afterword}
